
\begin{abstract}
Ulmer proposed to view a bialgebraic structure on an object $x$ as a structure map $f\colon Fx\to Gx$ satisfying an equation $\alpha_{(x,f)} = \beta_{(x,f)}$ of natural transformations. This perspective is powerful enough to account for many algebraic and coalgebraic structures, such as varieties of universal algebra, categories of algebras over a PROP, an operad or a monad. We develop Ulmer's bialgebras within a $2$-categorical setting with PIE limits, where structure maps are abstracted by inserters and their equations by equifiers. A syntax of signatures, theories and models is then used to produce these bialgebraic structures. The syntax of the commutative diagrams defining Hopf algebras can be used, as is, to define the $0$-cell of Hopf algebras within any symmetric pseudomonoid in a $2$-category with PIE limits.

In this setting, we show that lax morphisms of models induce morphisms between the objects of algebras, extending the classical result lifting a lax monoidal functor to a functor between the categories of internal monoids. Since the construction of the object of bialgebras is expressed entirely in terms of PIE limits, we show that Kan extensions, accessibility, local presentability, orthogonal factorization systems, regularity, and Barr-exactness lift along the construction under suitable assumptions. A bi-dense class of objects allows us to transfer monadicity and local presentability results to formal settings. Finally, we construct a bi-dense class of $\V$-categories with at most three objects, where $\V$ is any monoidal category with an absorbing initial object.
\end{abstract}

\smallskip
\noindent\textbf{Keywords:} Bialgebra, PIE limit, Algebraic theory, Monadicity, Locally presentable category, Exact category, Orthogonal factorization system

\smallskip
\noindent\textbf{MSC 2020:} 18C10, 18N10; 18A32, 18C15, 18C35, 18C40, 18E08, 16T05
